\section{Discussion}
\label{sec:discussion}

This study began with a diagnostic question and ended with a constrained method. The diagnostic stage showed that point performance, probability calibration, chemical similarity, conformal coverage, set informativeness, and selective risk were not interchangeable measures of trust. The design stage then used those disagreements as requirements rather than averaging them away. Under the locked Tox21 evaluation, the resulting TAME envelope improved the weakest class coverage by 1.36 percentage points on average, with a hierarchical-bootstrap interval above zero, while the 1.61-point MacroCSY loss remained inside the frozen 5-point non-inferiority bound. Because final predictions were sealed before labels were opened, this is independent evidence for the prespecified candidate in this assay family rather than a retrospective selection from favorable endpoint--method combinations.

\subsection{From a reliability illusion to an auditable intervention}

The ClinTox audit provides the clearest example of an aggregate reliability illusion. Marginal and shift-weighted conformal prediction appeared close to the nominal 0.90 target overall while positive-class coverage remained between 0.068 and 0.158. Mondrian calibration repaired the class deficit but returned both labels for approximately 71\% of compounds. Neither result is adequately summarized by \emph{coverage near 90\%}. The first is precise-looking but strongly asymmetric; the second is class-valid in the observed sample but often non-decisive. This distinction follows conformal work separating marginal from subgroup or class-conditional behavior and coverage from precision \citep{krstajic2021critical,cauchois2021knowing,bastani2022multivalid,ding2023classconditional,fisch2022trading}.

TAME does not attempt to erase that trade-off. It makes a one-sided, testable intervention: retain every ordinary Mondrian label, and add the development-identified protected label only when two audited transport views agree. The resulting invariants are stronger than a favorable empirical average. On any fixed prediction rows, class coverage cannot decrease, wrong-singleton exposure cannot increase, no empty set is emitted, and a new confident singleton cannot appear. What remains empirical is whether enough relevant baseline singletons are converted to deferrals on a new target cohort to produce a useful weakest-class gain without an unacceptable loss of correct singletons. The sealed EPA evaluation directly tested that remaining question.

\subsection{Transport diagnostics were decision gates, not decorative summaries}

Weighted conformal prediction is theoretically motivated when the relevant likelihood ratio is known \citep{tibshirani2019covariate}. In molecular applications the ratio is estimated in a selected representation, and that estimation step can dominate the result. In public development, the 2048-bit ECFP view produced complete coverage only by returning $\{0,1\}$ for every target compound; it also failed its effective-sample-size and clipping certificate in every primary cell. In the final cohort it again failed all 60 endpoint--seed audits. Relaxing thresholds after seeing this behavior would have converted an observed failure into an unrecorded method change.

The lower-dimensional physicochemical and score views behaved differently. Their inputs encode complementary notions of shift: changes in interpretable molecular properties and changes in the fitted ensemble's response geometry. Both passed in every final cell, while the public development physicochemical view failed in six cells and triggered exact fallback. This matters because a transport audit should be allowed to say \emph{do not transport}. Domain-classifier AUC alone is insufficient; the effective sample size, boundary clipping, class support, and whether weighting actually improves multivariate balance all affect whether a finite-sample correction is usable. Molecular implementations such as CoDrug show the promise of density-aware conformal prediction \citep{laghuvarapu2023codrug}; our results add evidence that a shift correction should expose its failure conditions rather than be applied automatically.

\subsection{Efficiency loss represents explicit deferral}

MacroCSY was chosen because ordinary set size alone does not distinguish a useful singleton from a wrong one. For a true class, correct-singleton yield asks how often the method returns exactly the correct label; averaging equally across classes prevents prevalence from dominating this efficiency measure. TAME can only preserve a singleton or turn it into a two-label set, so its MacroCSY cannot improve by construction. The scientific question was therefore non-inferiority, with the margin fixed before independent labels.

The final 1.61-point loss was well inside the 5-point bound, and every primary endpoint mean lay on the non-inferior side of that boundary descriptively. This should be read as a measured cost, not as a free gain. TAME improved coverage by converting some confident outputs into explicit deferrals. Whether that trade is desirable in a particular screening program depends on endpoint-specific costs of false exclusion, false inclusion, ambiguity, and experimental follow-up. The present study did not assign clinical or regulatory utility to those outcomes.

\subsection{Chemical similarity remains useful but cannot serve as a universal trust score}

The Stage~I applicability-domain results explain why ECFP novelty was retained as a diagnostic rather than promoted to a universal gate. Similarity ranked risk and miscoverage for BBBP and especially Lipophilicity, was weak for ClinTox, and reversed for ESOL under scaffold shift. A molecule can be structurally novel in one representation yet predictable because the endpoint depends on features shared across scaffolds; a close analogue can also be difficult in a steep or noisy local structure--property region. This endpoint dependence agrees with recent molecular uncertainty benchmarks and distance-based analyses \citep{lanini2024unique,kim2025distance,koetter2025upgrading,parrondo2026shifts}.

Selective prediction supplied a related warning. Confidence-based ClinTox rejection lowered ordinary risk partly by retaining only 27--31\% of positives when half of all observations remained. Risk--coverage curves can therefore combine error ranking with a change in class composition. Molecular abstention studies should report balanced risk and class-specific retention, while applicability-domain studies should report continuous associations and threshold sensitivity. A fixed Tanimoto threshold is a representation-specific operating choice, not a chemical constant.
This composition-aware interpretation is consistent with feature-aware predictive-inference work, which treats accepted-set behavior as part of the estimand rather than an incidental consequence of rejection \citep{teng2023feature}.

\subsection{Relation to contemporary ADMET model benchmarks}

The present contribution is complementary to architecture-centered ADMET studies. Broad benchmarks map the predictive capabilities of foundation models, graph neural networks, AutoML systems, and classical algorithms under realistic data challenges \citep{zhao2026revisiting}. Our representation-controlled audit instead isolates conflicts among post-prediction reliability quantities; RACER-C4 then tests a fixed uncertainty layer with an independent-label firewall. The design cannot establish where TAME would rank when attached to another representation, but it can expose a failure mode that a point-performance leaderboard does not measure.

This separation is important for reproducibility. Stage~I model families were never treated as independent replicates. Stage~II primary endpoints were selected by a public label-count rule, not by final performance. Development and prospective seeds were disjoint. Null and adverse endpoint results were retained, including the zero SR-ATAD5 minimum-class change and every negative MacroCSY change. Finally, the deterministic bootstrap repair changed only traversal order in an inference-only calculation; predictions, labels, point estimates, and the frozen interpretation were unchanged and both reports remain versioned.

\subsection{Limitations and scope}

The diagnostic panel contained four public endpoints, an ECFP4 representation, and conventional model families. It was designed to expose contrasting failure modes, not to represent the full ADMET task universe. The independent evaluation was restricted to the Tox21 nuclear-receptor and stress-response assay family and to one released external cohort. It was not a prospective wet-laboratory experiment, a clinical validation, or evidence of regulatory utility. The primary interval is conditional on the six prespecified endpoints, five seeds, locked preprocessing, and observed final cohort.

TAME also has deliberate structural limits. It only adds a development-selected label, so it cannot improve correct-singleton yield and will not repair a deficit in an unprotected label. The choice of label~0 was specific to the public development batch and must not be transferred to another assay family without a new, prospectively separated development process. The two transport audits are empirical diagnostics rather than proofs that density-ratio assumptions hold. Passing them does not supply conditional coverage or exact validity under arbitrary shift; failing them only certifies that the frozen augmentation is disabled.

Some external endpoint rows were excluded because standardized structures overlapped the training cohort, and two structures could not be parsed under the locked standardizer. Those exclusions were defined before labels and retained by identity, but they reduce the effective external sample. Rare classes also limit the resolution of endpoint-specific coverage estimates. Broader evaluation should include new assay families, temporal cohorts, learned molecular representations, and prospectively generated experimental data while preserving the same separation between design and final labels.
Assay labels can also contain biological and measurement heterogeneity that no structural split or conformal wrapper removes; robustness of marginal coverage to label noise does not by itself resolve the scientific meaning of the target \citep{einbinder2023labelnoise}.

\subsection{Conclusion}

Reliable molecular prediction cannot be reduced to aggregate accuracy, marginal coverage, or a single distance-to-training score. The four-endpoint audit showed how each can fail: structured splits changed performance, marginal coverage concealed class undercoverage, class-conditional repair could become ambiguous, similarity was endpoint dependent, and selective rejection could change who remained. TAME converted those observations into an auditable intervention with explicit transport certificates, exact fallback, baseline containment, and a prediction-to-label firewall. In the locked EPA evaluation it improved minimum-class coverage while remaining within the prespecified MacroCSY efficiency bound. The defensible conclusion is narrow but useful: failure-aware constraints and sealed validation can improve the credibility of a molecular uncertainty layer, while generalization beyond the evaluated Tox21 setting remains an empirical question.
